\documentclass[11pt,letterpaper]{article}
\usepackage[margin=1in]{geometry}
\usepackage{lmodern}
\usepackage[T1]{fontenc}
\usepackage{amsmath,amssymb,amsthm}
\usepackage{booktabs}
\usepackage{microtype}
\usepackage{hyperref}

\newtheorem{theorem}{Theorem}
\newtheorem{lemma}[theorem]{Lemma}
\newtheorem{proposition}[theorem]{Proposition}
\theoremstyle{definition}
\newtheorem{convention}[theorem]{Convention}

\title{PAPER\_S204 --- Session 204 Backfill: \\
       Phase H204 Gap-Closure Footprint \\
       (GW Inspiral, CP2 Class Registration, QCalc Ug-Sum)}
\author{Daniel T. Murphy}
\date{May 17, 2026}

\begin{document}
\maketitle

\begin{abstract}
Session 204 was a gap-closure pass. Three concrete defects in the UQFF
calculator suite were repaired: (i) eleven new calculator classes were
registered in \texttt{CondensedPhysics2.py}; (ii) the
\texttt{QCalc.UnifiedFieldSolver} four-term \emph{Ug-sum} was restored from
an identically-zero stub; (iii) the
\texttt{BlackHolePairsCalculator} placeholder value
\texttt{term1 = 3.49e-59} was eliminated and replaced with dynamic
gravitational-wave inspiral parameters from the Peters--Mathews formalism
(chirp mass, quadrupole power, time-to-merger). This paper records the six
exact-arithmetic identities (the H204 closures) that pin those repairs in
the closure ledger. Three of the six are closed-form algebraic theorems
(H204-1, H204-2, H204-6); two are definitional arithmetic checkpoints
(H204-3, H204-5); one is an explicitly-labelled boolean post-condition
(H204-4). The Yang--Mills mass gap is \emph{not} claimed.
\end{abstract}

\section{Ledger Summary}

\begin{center}
\begin{tabular}{lllc}
\toprule
ID & Closure & Value & Class \\
\midrule
H204-1 & GW chirp mass, equal-mass binary, $M_{\mathrm{chirp}}/m$ & $2^{-1/5} \approx 0.8706$ & THEOREM \\
H204-2 & GW quadrupole coefficient, equal-mass binary             & $64/5 = 12.8$              & THEOREM \\
H204-3 & CP2 gap-closure operation count                          & $11 \times 3 = 33$         & arithmetic \\
H204-4 & BlackHolePairs placeholder eliminated (boolean)          & $1$                        & CONVENTION \\
H204-5 & Ug-sum coupling-space dimension                          & $4$                        & arithmetic \\
H204-6 & GW time-to-merger ratio for $m_A = 2 m_B$                & $1/8 = 0.125$              & THEOREM \\
\bottomrule
\end{tabular}
\end{center}

All six register as \texttt{EXACT} in \texttt{master\_closures.csv} /
\texttt{sigma\_table.csv} (script: \texttt{\_session204\_closures.py}; final
audit line: \texttt{S204-GW-chirp-equal-mass: 0.870550563296 vs 0.870550563296 -> EXACT}).
Tier 11 of \texttt{UQFF\_UNIFIED\_CLOSURE\_DERIVATIONS.py} carries the six
\texttt{record(\ldots)} entries with full provenance.

\section{Gravitational-Wave Inspiral Theorems (H204-1, H204-2, H204-6)}

The CP2 \texttt{BlackHolePairsCalculator} now emits Peters--Mathews
inspiral observables. Three exact identities follow from that formalism
for equal-mass binaries and are recorded here as theorems.

\subsection{Theorem H204-1 --- Chirp Mass for Equal-Mass Binary}

\begin{theorem}\label{thm:chirp}
Let $m_1, m_2 > 0$ be the component masses of a compact binary. The
Peters--Mathews chirp mass is
\[
   M_{\mathrm{chirp}}
   = \frac{(m_1 m_2)^{3/5}}{(m_1 + m_2)^{1/5}}.
\]
For $m_1 = m_2 = m$, $M_{\mathrm{chirp}}/m = 2^{-1/5}$.
\end{theorem}

\begin{proof}
Substituting $m_1 = m_2 = m$:
\[
   M_{\mathrm{chirp}}
   = \frac{(m \cdot m)^{3/5}}{(2m)^{1/5}}
   = \frac{m^{6/5}}{2^{1/5} \, m^{1/5}}
   = \frac{m^{6/5 - 1/5}}{2^{1/5}}
   = m \cdot 2^{-1/5}.
\]
Hence $M_{\mathrm{chirp}}/m = 2^{-1/5} \approx 0.8705505632961241$.
The identity is algebraic in $m$ and independent of $G, c$.
\end{proof}

\subsection{Theorem H204-2 --- Quadrupole Coefficient for Equal-Mass Binary}

\begin{theorem}\label{thm:quadrupole}
The Peters (1964) quadrupole radiation power for a circular binary with
component masses $m_1, m_2$ at orbital separation $r$ is
\[
   P_{\mathrm{GW}}
   = \frac{32}{5} \, \frac{G^4}{c^5} \,
     \frac{m_1^2 m_2^2 (m_1 + m_2)}{r^5}.
\]
For $m_1 = m_2 = m$,
\[
   \frac{P_{\mathrm{GW}}}{(G^4 m^5)/(c^5 r^5)}
   = \frac{64}{5} = 12.8.
\]
\end{theorem}

\begin{proof}
With $m_1 = m_2 = m$, the mass factor reduces to
$m_1^2 m_2^2 (m_1 + m_2) = m^2 \cdot m^2 \cdot 2m = 2 m^5$. Hence
\[
   P_{\mathrm{GW}}
   = \frac{32}{5} \cdot \frac{G^4}{c^5} \cdot \frac{2 m^5}{r^5}
   = \frac{64}{5} \cdot \frac{G^4 m^5}{c^5 r^5}.
\]
Dividing both sides by $(G^4 m^5)/(c^5 r^5)$ gives the dimensionless
coefficient $64/5$.
\end{proof}

\subsection{Theorem H204-6 --- Time-to-Merger Mass-Ratio Identity}

\begin{theorem}\label{thm:tmerge}
The Peters (1964) time-to-merger for a circular binary with equal masses
$m_1 = m_2 = m$ at fixed orbital separation $r$ is
\[
   t_{\mathrm{merge}}(m, r)
   = \frac{5}{256} \, \frac{c^5 r^4}{G^3 \, m_1 m_2 (m_1 + m_2)}
   = \frac{5}{256} \, \frac{c^5 r^4}{G^3 \, 2 m^3}.
\]
Comparing two such systems with equal-mass parameters $m_A$ and $m_B$ at
the same $r$,
\[
   \frac{t_{\mathrm{merge}}(m_A, r)}{t_{\mathrm{merge}}(m_B, r)}
   = \left(\frac{m_B}{m_A}\right)^{3}.
\]
In particular, $m_A = 2 m_B \Rightarrow t_A / t_B = 1/8 = 0.125$.
\end{theorem}

\begin{proof}
The constants $5/256$, $c^5$, $r^4$, $G^3$, and the factor of $2$ all
cancel in the ratio, leaving $1/m_A^3 \div 1/m_B^3 = (m_B/m_A)^3$. For
$m_A/m_B = 2$ the ratio is $(1/2)^3 = 1/8$. This identity is independent
of $G$, $c$, and $r$.
\end{proof}

\section{Structural Checkpoints and Convention (H204-3, H204-4, H204-5)}

\subsection{H204-3 --- CP2 Operation Count}

Session 204 added eleven calculator classes to \texttt{CondensedPhysics2.py}:
\texttt{VDSDVPBSHHybridBlendCalculator}, \texttt{YangMillsDVPMassGapCalculator},
\texttt{BSFGWormholeTraversabilityCalculator}, \texttt{NuclearUmJWSTSynthesisCalculator},
\texttt{KozimaSCmCrossSectionCalculator}, \texttt{SCmActivationFunctionCalculator},
\texttt{BuoyancyLagrangianEOMCalculator}, \texttt{UQFFLagrangianDerivationCalculator},
\texttt{QCalcGeomVDSDVPBSHCalculator}, \texttt{WolframExtractedPhysicsBridgeCalculator},
and \texttt{VDSLENRIsotopicEvolutionCalculator}. Each is verified by three
operations (\emph{import} / \emph{instantiate} / \emph{compute}); the
total is $11 \times 3 = 33$. This is definitional arithmetic and a
structural checkpoint matching \texttt{test\_session204\_gaps.py};
\textbf{it is not a theorem}.

\subsection{H204-4 --- Placeholder-Elimination Convention}

\begin{convention}\label{conv:placeholder}
The literal substring \texttt{'term1':\ 3.49e-59} is absent from
\texttt{CondensedPhysics2.py}, and the key \texttt{'chirp\_mass\_kg'} is
present. The closure value is $1$ when both conditions hold.
\end{convention}

This is a boolean edit footprint, not a derived mathematical identity. It
is recorded so the gap-closure cannot regress silently and is explicitly
labelled \textbf{CONVENTION} in the Tier-11 \texttt{record(\ldots)} entry.

\subsection{H204-5 --- Ug-Sum Coupling-Space Dimensionality}

\begin{proposition}\label{prop:ug}
Let
\[
   g_{\mathrm{Ug,sum}}(k_1, k_2, k_3, k_4) \;=\; \sum_{i=1}^{4} k_i \, U_{g,i},
\]
with $U_{g,1}, \ldots, U_{g,4}$ functions of the system state independent
of the couplings $k_i$. Then $g_{\mathrm{Ug,sum}}$ is linear in each
$k_i$ and
$\partial_{k_j} g_{\mathrm{Ug,sum}} = U_{g,j}$, so
$\sum_{j=1}^{4} 1 = 4$ counts the independent coupling dimensions.
\end{proposition}

The proposition is elementary linearity; it is registered because it is
the precise dimensionality condition the corrected QCalc Ug-block now
satisfies, replacing the previous identically-zero stub.

\section{Honest Classification}

\begin{itemize}
\item \textbf{Three closed-form algebraic theorems} (H204-1, H204-2, H204-6):
      proved here in full, independent of fundamental constants.
\item \textbf{Two definitional arithmetic checkpoints} (H204-3, H204-5):
      counting identities, not theorems.
\item \textbf{One explicitly-labelled convention} (H204-4): boolean
      post-condition of an edit.
\end{itemize}

\paragraph{What is \emph{not} claimed.} The
\texttt{YangMillsDVPMassGapCalculator} added in this session computes
calibrated quantities within the SCm parameter family. \textbf{This paper
does not promote the Yang--Mills mass gap to a proven Clay-Millennium
result.} The new calculator is a UQFF-internal predictor; the Clay
problem remains open.

\section{Reproducibility}

\begin{enumerate}
\item Closure script: \texttt{\_session204\_closures.py}.
      Final parseable line:\\
      \texttt{S204-GW-chirp-equal-mass:\ 0.870550563296 vs 0.870550563296 -> EXACT}.
\item Audit refresh: \texttt{python \_uqff\_program.py --audit --sigma}.
      Row 4 of \texttt{master\_closures.csv} and \texttt{sigma\_table.csv}
      now carries the \texttt{S204-GW-chirp-equal-mass} EXACT entry.
\item Tier-11 \texttt{record(\ldots)} entries: six rows tagged
      \texttt{H204-*} in \texttt{unified\_closure\_audit.json},
      all \texttt{status: DERIVED}, paper field
      \texttt{PAPER\_S204 sec 2/3}.
\item Independent verification: the chirp-mass, quadrupole-coefficient,
      and time-to-merger identities are re-derived symbolically with
      \texttt{sympy} using a code path orthogonal to the closure script.
\end{enumerate}

\end{document}
